\input format.tex

%\usepackage{arydshln}

\begin{document}

%%%
\setlength{\arrayrulewidth}{0.5pt}
\fontsize{9.3pt}{17pt}\selectfont
\color{gray2}

\vspace*{0mm}

\bigskip
\noindent
基于本次您的检测结果，建议您从以下几个方面进行干预：
\begin{compactenum}[\mbox{1、}]
\xiaowuhao \color{gray2} \lantxh {\item[1、]
鉴于您的肠道菌群多样性处于正常范围，请继续保持良好的饮食及生活习惯，不滥用抗生素，可适当多喝酸奶，帮助维持肠道菌群平衡；您体内未超标的致病菌在您免疫能力低下时，可能诱发相关疾病，建议您持续监测肠道健康，必要时到正规医院进行相关检查。
}
\xiaowuhao \color{gray2} \lantxh {\item[2、]
您的肠道菌群不会增加您患2型糖尿病、结直肠癌、炎症性肠病的风险，请您继续保持良好的饮食及生活习惯，维护肠道健康。
}
\xiaowuhao \color{gray2} \lantxh {\item[3、]
您的肠道菌群检测结果显示您的消化吸收及相关代谢情况良好，请您继续保持，避免滥用抗生素等导致肠道菌群失调。
}
\end{compactenum}

\bigskip
\tabcolsep=6pt
\fontsize{9.3pt}{9.8pt}\selectfont
\begin{tctabularx}{tabularx={p{15.4cm}<{\centering}}}
\\[-6.769pt]
\cellcolor{kelly} \raisebox{1pt}{\color{black70}\xiaosihao\lantxh\bf 膳食方案} \\
\end{tctabularx}

{
\lantxh
\vspace*{-4.25mm}
\noindent
\begin{longtable}{|p{1.25cm}@{}p{2cm}|p{11.732cm}|}
\hline
\parbox[c]{\hsize}{\vskip7pt\raisebox{-1.5ex}{\includegraphics[width=.9cm]{zhushi.pdf}}\vskip7pt} & \parbox[c]{\hsize}{\centering\xiaowuhao \color{gray2}\vskip7pt主食\vskip7pt} & \parbox[c]{\hsize}{\xiaowuhao \color{gray2}\vskip7pt减少主食摄入量，粗细粮搭配；粗粮最好占主食的一半。\vskip7pt}\\
\hline
\parbox[c]{\hsize}{\vskip7pt\raisebox{-1.5ex}{\includegraphics[width=.9cm]{roulei.pdf}}\vskip7pt} & \parbox[c]{\hsize}{\centering\xiaowuhao \color{gray2}\vskip7pt肉类\vskip7pt} & \parbox[c]{\hsize}{\xiaowuhao \color{gray2}\vskip7pt宜吃禽类，如去皮鸡肉；避免摄入肥猪肉、肥牛、猪脑等高脂食物，适量减少五花肉、排骨的摄入。\vskip7pt}\\
\hline
\parbox[c]{\hsize}{\vskip7pt\raisebox{-1.5ex}{\includegraphics[width=.9cm]{yuxialei.pdf}}\vskip7pt} & \parbox[c]{\hsize}{\centering\xiaowuhao \color{gray2}\vskip7pt鱼虾类\vskip7pt} & \parbox[c]{\hsize}{\xiaowuhao \color{gray2}\vskip7pt适量多食用鱼虾类，代替畜肉类。\vskip7pt}\\
\hline
\parbox[c]{\hsize}{\vskip7pt\raisebox{-1.5ex}{\includegraphics[width=.9cm]{danlei.pdf}}\vskip7pt} & \parbox[c]{\hsize}{\centering\xiaowuhao \color{gray2}\vskip7pt蛋奶、大豆类\vskip7pt} & \parbox[c]{\hsize}{\xiaowuhao \color{gray2}\vskip7pt蛋类每天1个，牛奶宜选用脱脂奶，适量食用豆制品，如豆腐、豆浆等。\vskip7pt}\\
\hline
\parbox[c]{\hsize}{\vskip7pt\raisebox{-1.5ex}{\includegraphics[width=.9cm]{shucaishuiguo.pdf}}\vskip7pt} & \parbox[c]{\hsize}{\centering\xiaowuhao \color{gray2}\vskip7pt蔬菜水果\vskip7pt} & \parbox[c]{\hsize}{\xiaowuhao \color{gray2}\vskip7pt多食用新鲜蔬菜，其中深色蔬菜占1{/}2，多吃富含维生素的蔬菜，如芹菜、韭菜等；少吃含糖量高的水果，如枣、龙眼、荔枝、西瓜等。\vskip7pt}\\
\hline
\parbox[c]{\hsize}{\vskip7pt\raisebox{-1.5ex}{\includegraphics[width=.9cm]{you.pdf}}\vskip7pt} & \parbox[c]{\hsize}{\centering\xiaowuhao \color{gray2}\vskip7pt食用油、盐\vskip7pt} & \parbox[c]{\hsize}{\xiaowuhao \color{gray2}\vskip7pt食用油不超过25g，少吃动物油，食盐不超过6g。\vskip7pt}\\
\hline
\parbox[c]{\hsize}{\vskip7pt\raisebox{-1.5ex}{\includegraphics[width=.9cm]{other.pdf}}\vskip7pt} & \parbox[c]{\hsize}{\centering\xiaowuhao \color{gray2}\vskip7pt其他\vskip7pt} & \parbox[c]{\hsize}{\xiaowuhao \color{gray2}\vskip7pt控制总能量的摄入，每天总能量(kcal)=理想的体重(kg)*(20\textasciitilde 25)(kcal{/}kg) 。全天能量分配：早餐30{\%}、午餐40{\%}、晚餐30{\%}。改变不良饮食习惯和行为，戒烟酒。烹调方法应尽量选择拌，炖、蒸、焖方法，忌煎、炸、烧、烤、熏等方法。不吃高热量食物，如油炸薯条、方便面、含糖饮料、油饼、油条等。如患痛风、溃疡性结肠炎、慢性腹泻、食物不耐受等有饮食禁忌的疾病，请综合考虑疾病的饮食原则。\vskip7pt}\\
\hline
\end{longtable}
\vspace*{-6mm}
\begin{longtable}{p{0.1cm}p{15.1cm}}
\includegraphics{xiaofangkuai.pdf} & {\qihao 以上膳食方案仅供参考，您的实际饮食请结合自身具体情况调整。} \\
\end{longtable}

}
\end{document}
